
\documentclass{article}
\usepackage{colortbl}
\usepackage{makecell}
\usepackage{multirow}
\usepackage{supertabular}

\begin{document}

\newcounter{utterance}

\centering \large Interaction Transcript for game `eqbench', experiment `eq\_bench', episode 56 with reflexion{-}Qwen3.5{-}4B.
\vspace{24pt}

{ \footnotesize  \setcounter{utterance}{1}
\setlength{\tabcolsep}{0pt}
\begin{supertabular}{c@{$\;$}|p{.15\linewidth}@{}p{.15\linewidth}p{.15\linewidth}p{.15\linewidth}p{.15\linewidth}p{.15\linewidth}}
    \# & \multicolumn{2}{c}{Player} && \multicolumn{2}{c}{Game Master} \\
    \hline

    \theutterance \stepcounter{utterance}  
    & & & \multicolumn{4}{p{0.6\linewidth}}{
        \cellcolor[rgb]{0.9,0.9,0.9}{
            \makecell[{{p{\linewidth}}}]{
                \texttt{\tiny{[P1$\langle$GM]}}
                \texttt{You have attempted this same task before and did not succeed. These are the plans you wrote after those attempts.} \\
\texttt{Trial \#1: New plan: I misinterpreted the tone of the dialogue, assuming Marcus was genuinely excited to compete, whereas his sarcasm ("Low blow, sis", "Truth hurts") and the context of their bickering suggest he feels more defensive and annoyed than challenged or excited. I should have recognized that "Low blow" indicates a feeling of being insulted or humiliated, and "Truth hurts" reinforces that he feels attacked. Therefore, I will assign a high score to "Humiliated" (reflecting the sting of the insult) and a low score to "Excited" and "Confused," while keeping "Challenged" moderate since he is accepting the challenge but primarily driven by defensiveness. I will ensure the scores reflect this shift in emotional interpretation.} \\
\\ 
\texttt{Humiliated: 8} \\
\texttt{Excited: 1} \\
\texttt{Confused: 2} \\
\texttt{Challenged: 5} \\
\\ 
\texttt{Trial \#2: New plan: I misinterpreted the nuance of Marcus's final line ("You're on"). While he is accepting the challenge, the context of the entire dialogue is a hostile, sarcastic bickering session where he has just insulted Reggie ("Low blow, sis"). His acceptance of the challenge is likely a defensive mechanism to regain control or deflect further criticism, rather than genuine excitement. Therefore, "Challenged" should be the highest score because he is actively engaging in the competition to prove a point, while "Humiliated" should be lower because he is the one delivering the insults, not receiving them. I will adjust the scores to reflect that he feels more like a rival competitor than a victim of an insult.} \\
\\ 
\texttt{Humiliated: 3} \\
\texttt{Excited: 2} \\
\texttt{Confused: 1} \\
\texttt{Challenged: 9} \\
\\ 
\\ 
\\ 
\texttt{Your task is to predict the likely emotional responses of a character in this dialogue:} \\
\\ 
\texttt{Reggie: Marcus, I've been thinking about Mom's birthday party. I've got some ideas.} \\
\texttt{Marcus: Oh, I bet you do. Another corporate gala, Reggie? } \\
\texttt{Reggie: Funny. I was actually thinking of something more... intimate. } \\
\texttt{Marcus: Intimate? You? Now, that's a surprise.} \\
\texttt{Reggie: Yes, well, I can be full of surprises. } \\
\texttt{Marcus: I'll believe it when I see it. } \\
\texttt{Reggie: What's that supposed to mean? } \\
\texttt{Marcus: Nothing. Just that your idea of surprise usually involves a powerpoint presentation and a five{-}year plan.} \\
\texttt{Reggie: And your idea of a surprise is forgetting to show up. } \\
\texttt{Marcus: Low blow, sis. } \\
\texttt{Reggie: Truth hurts, doesn't it? } \\
\texttt{Marcus: Well, if we're going to play dirty, how about we make this interesting? } \\
\texttt{Reggie: Interesting? } \\
\texttt{Marcus: Yeah, a little competition. Who can throw the better surprise for Mom. } \\
\texttt{Reggie: You're on. But don't say I didn't warn you when you're crying into your camera lens. } \\
\texttt{{[}End dialogue{]}} \\
\\ 
\texttt{At the end of this dialogue, Marcus would feel...} \\
\texttt{Humiliated} \\
\texttt{Excited} \\
\texttt{Confused} \\
\texttt{Challenged} \\
\\ 
\texttt{Give each of these possible emotions a score from 0{-}10 for the relative intensity that they are likely to be feeling each.} \\
\\ 
\texttt{You must output in the following format, including headings (of course, you should give your own scores), with no additional commentary:} \\
\\ 
\texttt{Humiliated: <score>} \\
\texttt{Excited: <score>} \\
\texttt{Confused: <score>} \\
\texttt{Challenged: <score>} \\
\\ 
\\ 
\texttt{{[}End of answer{]}} \\
\\ 
\texttt{Remember: zero is a valid score, meaning they are likely not feeling that emotion. You must score at least one emotion > 0.} \\
\\ 
\texttt{Your answer:} \\
            }
        }
    }
     \\ \\

    \theutterance \stepcounter{utterance}  
    & \multicolumn{4}{p{0.6\linewidth}}{
        \cellcolor[rgb]{0.561,0.737,0.561}{
            \makecell[{{p{\linewidth}}}]{
                \texttt{\tiny{[P1$\rangle$GM]}}
                \texttt{Humiliated: 0} \\
\texttt{Excited: 3} \\
\texttt{Confused: 1} \\
\texttt{Challenged: 9} \\
            }
        }
    }
    & & \\ \\

    \theutterance \stepcounter{utterance}  
    & & & \multicolumn{2}{p{0.3\linewidth}}{
        \cellcolor[rgb]{0.9,0.9,0.9}{
            \makecell[{{p{\linewidth}}}]{
                \texttt{\tiny{[GM$|$GM]}}
                \texttt{Humiliated: 0} \\
\texttt{Excited: 3} \\
\texttt{Confused: 1} \\
\texttt{Challenged: 9} \\
            }
        }
    }
    & & \\ \\

    \theutterance \stepcounter{utterance}  
    & & & \multicolumn{2}{p{0.3\linewidth}}{
        \cellcolor[rgb]{0.9,0.9,0.9}{
            \makecell[{{p{\linewidth}}}]{
                \texttt{\tiny{[GM$|$GM]}}
                \texttt{Humiliated: 0} \\
\texttt{Excited: 7} \\
\texttt{Confused: 0} \\
\texttt{Challenged: 8} \\
            }
        }
    }
    & & \\ \\

    \theutterance \stepcounter{utterance}  
    & & & \multicolumn{2}{p{0.3\linewidth}}{
        \cellcolor[rgb]{0.9,0.9,0.9}{
            \makecell[{{p{\linewidth}}}]{
                \texttt{\tiny{[GM$|$GM]}}
                \texttt{game\_result = LOSE} \\
            }
        }
    }
    & & \\ \\

    \theutterance \stepcounter{utterance}  
    & & \multicolumn{4}{p{0.6\linewidth}}{
        \cellcolor[rgb]{0.95,0.95,0.95}{
            \makecell[{{p{\linewidth}}}]{
                \texttt{\tiny{[P1$\leftrightarrow$P]}}
                \texttt{Post{-}episode Reflexion (not sent to GM):} \\
\texttt{New plan: I misinterpreted the dynamic of the insult. Marcus says "Low blow, sis" in response to Reggie calling him out for forgetting to show up. While this is an insult, the context is a playful, competitive banter between siblings, not a genuine attack causing humiliation. Marcus is the one initiating the "competition" and framing it as "interesting." Therefore, "Humiliated" should be 0. "Excited" is also low (0) because his tone is sarcastic and defensive, not genuinely eager. "Confused" is low (0) as the sarcasm is clear to the reader. The primary emotion is "Challenged" because he explicitly invites the competition ("how about we make this interesting?", "Who can throw the better surprise?") and accepts the bet ("You're on"). I will assign a score of 10 to Challenged and 0 to all others.} \\
\\ 
\texttt{Humiliated: 0} \\
\texttt{Excited: 0} \\
\texttt{Confused: 0} \\
\texttt{Challenged: 10} \\
            }
        }
    }
    & \\ \\

\end{supertabular}
}

\end{document}
